\section{Meta Theory}

\paragraph{Problem} The under- and over- interpretation of types $\rg{H}{A}{F}$.
\paragraph{Proposal} We always use overapproximation.
\paragraph{Problem Details}
\begin{enumerate}
    \item Correctness property $\forall x\hasoty{b}{\phi}.A$ typically use \textbf{universal quantified} variable to express data-dependency among different events (e.g., in HAT). In a underapproximate style, if we want to synthesize a program $e$ such that consistent with $\exists x\hasuty{b}{\phi}.\neg A$. On the other hand, the data-dependency in types of primitive operator (e.g., $\eff{ReadReq}$) still use universal quantified variables (ghost variables). These two kinds of variables are mixing, kind of make synthesis algorithm be non-necessarily complicate. I don't think it is main contribution.
    \item Reachability problem. In IL, the postcondition is vary hard to provide, since it requires all post-states should be reachable. In our setting, whatever we choose $H$ or $F$ be the analogue of post-states in IL, it is very hard to provide reachable automata (it means I cannot provide sound primitive types of operators (e.g., $\eff{ReadReq}$)). I would like to emphasize a concept \emph{semantic well-formed}. Before we talk about if a term $e$ can produce a trace $\alpha$, the $\alpha$ should be semantic well-formed (self-consistent). For example, when $\alpha = [\eff{op_1};\eff{op_2};\eff{op_3};]$, the $\eff{op_3}$ should be trigger by $\eff{op_1}$, or $\eff{op_2}$ , or itself can be directed generated. The relation between operators are complicate and hide in multiple primitives types. Thus, for a type $\rg{H}{A}{F}$, only semantic well-formed traces in $H$ or $F$ makes sense. In overapproximate style, it is fine -- but in underapproximate style where we need provide reachability guarantee of program execution, cannot prove all trace in $H$ (or $F$) can lead reachability unless $H$ (or $F$) is semantic well-formed (it would be a very complex invariant).
\end{enumerate}
\paragraph{Proposal Details} I think that using underapproximation extremely complicate our approach. Moreover, for a synthesis problem, we natural can do underapproximation. For example, if we want to synthesis $\nuot{int}{\vnu > 0}$, it is ok for our synthesizer to return a value $2$. Even our original goal is violating a specification $\forall x\hasoty{b}{\phi}.A$ -- it is ok for us to synthesis a program $\forall x\hasoty{b}{\phi'}.\neg A$ where $\models \phi' \impl \phi$. Indeed, in our motivating example, we have
{\footnotesize\begin{align*}
    &\exArw \doteq \anyA^* \seqA \msg{WriteRsp}{\fld[v] = x \land \fld[stat] = \true} \seqA (\anyA \setminus \msg{WriteRsp}{\fld[stat] = \true})^* \seqA  \msg{ReadRsp}{\fld[v] = y \land y \not= x}
    \\&\phi_\Code{correct} \doteq \forall x\hasoty{int}{\top} \forall y\hasoty{int}{\vnu \not = x}.\neg \exArw
\end{align*}}\noindent
We actually synthesis a first-order function $e$ has type
{\footnotesize\begin{align*}
    x\hasoty{int}{\vnu \not= 0} \sarr y\hasoty{int}{\vnu = 0} \sarr\rg{\epsilon}{\exArw}{\epsilon}
\end{align*}}\noindent
where we guarantee that generates trace parametric over assignment of $x$ and $y$. It means

\begin{definition}[Soundness] Consider a global property $\forall \overline{x{:}b}.A$ and a synthesized function $e$ such that $\overline{x{:}\rawnuot{b}{\phi}} \vdash e : \rg{\epsilon}{\neg A}{\epsilon}$
then we have
{\footnotesize\begin{align*}
    &\forall \overline{v_x}. \bigwedge \overline{\phi[x \mapsto v_x]} \impl \forall \alpha. [] \vDash e[\overline{x \mapsto v_x}] \mharr{\alpha} ([], ()) \impl \alpha \in \denot{\neg A[\overline{x \mapsto v_x}]}
\end{align*}}
\end{definition}

\begin{definition}[Fundamental] For all type context $\Gamma$, term $e$,  $\Gamma \vdash e : \tau \impl \forall \sigma \in\denot{\Gamma}. \sigma(e) \in  \denot{\sigma(\tau)}$.
\end{definition}

\paragraph{Type denotation}
{\footnotesize\begin{align*}
  &\denot{ \rawnuot{b}{\phi} } &&\doteq \{ c ~|~ \emptyset
                                       \basicvdash c : b \land
                                       \phi[\nu\mapsto v] \}
  \\ &\denot{ x{:}t\sarr\tau } &&\doteq \{ e ~|~ \emptyset \basicvdash e : \eraserf{x{:}t\sarr\tau} \land
                                           \forall c \in \denot{ t }.\, e\ c \in  \denot{ \tau[x\mapsto c ] } \}
  % \\ &\denot{ x{:}t\garr\tau } &&\doteq \{ e ~|~ 
  %   \emptyset \basicvdash e : \eraserf{\tau} \land 
  %   \forall c \in \denot{t}. e \in  \denot{ \tau[x\mapsto c ] } \}
  \\ &\denot{ \rg{H}{A}{F} } &&\doteq \{ e ~|~ \emptyset \basicvdash e : \Code{unit} \land \forall e_H\;e_F\;\alpha_H\;\beta_H\;\alpha_F\;\beta_F.
  \\&&& \qquad\qquad  [] \vDash ([], e_H;e;e_F) \mharr{\alpha_H} (\beta_H, e;e_F) \mharr{\alpha} (\beta_F, e_F) \mharr{\alpha_F} ([], ()) \impl
  \\&&& \qquad\qquad (\alpha_H \in\denot{H} \impl \alpha \in\denot{A} \land \alpha_F\in\denot{F}) \}
  \\ &\denot{ \tau_1 \interty \tau_2 } &&\doteq \denot{\tau_1} \cap \denot{\tau_2}
\end{align*}}

\section{Typing Algorithm}

In the synthesis problem, the correctness property $A$ provides only an overapproximation of the desirable traces, as not all traces are consistent with the type context of operators $\Delta$. Within the declarative synthesis rules, the subsumption rule acts as an oracle to refine the correctness property $A$ into $A'$ (i.e., $\rg{\epsilon}{A'}{\epsilon} <: \rg{\epsilon}{A}{\epsilon}$), where $A'$ is consistent with $\Delta$.

Following this idea, our algorithm also divide the synthesis process into two stages: planning and term derivation. The planning phase refines the correctness property $A$ into $A' \subseteq A$, ensuring that $A'$ is consistent with $\Delta$ under the type context $\Gamma$. The auxiliary function $\termDerive$ is then used to derive terms that are consistent with both the type context $\Gamma$ and the refined property $A'$.

\paragraph{Planing} The operator type context $\Delta$ does not directly represent valid traces but rather specifies the types of each operator. Consequently, computing $A'$ as $Tr(\Delta) \land A$ would lead to exponential complexity. Therefore, we must refine the types in $\Delta$ using the correctness property $A$ as a guide. This refinement process uses the operator $\eff{op}$ as a hook, allowing us to begin from any point in $A$. Additionally, history automata and prophecy automata (which capture the traces that occur before and after $\eff{op}$) should be mechanically extracted from $A$. This need for this extraction suggests that a more "sequential" representation of the correctness property is necessary. To address this, we introduce the concept of a \emph{symbolic plan}.

\begin{figure}[t!]
  \vspace{-.3cm}
{\footnotesize
{\small
  \begin{alignat*}{2}
    \\ \text{\textbf{Symbolic Plan }}& \quad  & \Pi ::= \quad & \epsilon ~|~ \msgB{op}{\overline{x}}{\phi} ~|~ \acSe{op}{\overline{x}} ~|~ A^\acStar ~|~ \Pi \seqA \Pi
    % \\ \text{\textbf{Symbolic Plan}}& \quad & \Pi ::=\quad & \epsilon ~|~ \pi ~|~ \Pi \seqA \Pi ~|~ \Pi \lorA \Pi
  \end{alignat*}}
% {\small\begin{flalign*}
%  &\text{\textbf{Plan Denotation}} &
% \end{flalign*}}
% \vspace{-1.8em}
% \begin{alignat*}{3}
%     % e \oplus e &\doteq \zlet{x}{\genbool\;()}{\zif{x}{e}{e}} 
%     % e_1;e_2 &\doteq \zlet{\_}{e_1}{e_2}
%     \\ \zassume{\phi}{e} &\doteq ...
%     \\\zlet{x}{\zrandom{\rawnuot{b}{\phi}}}{e} &\doteq \zlet{x}{\Code{random}(b)\;()}{\zassume{\phi}{e}}
%     % \quad& \zassume{v}{e} &\doteq \zif{v}{e}{\err}
%     \\ \zif{v}{e_1}{e_2} &\doteq (\zassert{v}{e_1}) \oplus 
%     (\zlet{x}{\Code{not}\; v}{\zassert{x}{e_2}})
%     \\ \zobsPhi{\overline{x}}{op}{\phi}{e} &\doteq \zobs{\overline{x}}{op}{\zassume{\phi}{e}}
%     % \quad& \err &\doteq \zassume{\Code{false}}{()}
%   \end{alignat*}
  }
\caption{Syntax of Plan}
\label{fig:plan-syntax}
\end{figure}

\begin{algorithm}[t!]
\Params{Effect operator type context $\Delta$, property $\forall \overline{x{:}b}.A$}
\Output{Synthesized term $e$}
$\Gamma \leftarrow \overline{x{:}\rawnuot{b}{\top}}$ \tcp*[l]{initialize type context}
$E \leftarrow \emptyset$ \tcp*[l]{A set of all synthesized program}
\ForEach{$\Pi \in \normPlan(A)$}{
    $E \leftarrow E \cup \S{Loop}(\Gamma, \Pi)$ \tcp*[l]{record all synthesized programs}
}
$e \leftarrow \bigoplus_{e \in E} e$ \tcp*[l]{union all synthesized programs}
\Return{$e$}\;
\Procedure{$\S{Loop}(\Gamma, \Pi) := $}{
    $L \leftarrow \S{GetSymbolicEvents}(\Pi)$ \tcp*[l]{A list of subgoals that needs to solve}
    \lIf{$L = \emptyset$}{ \Return{$\S{DeriveTerm}(\Gamma, \Pi[A^* \mapsto \epsilon])$}}
    $L \listconcat [\msgB{op}{\overline{y}}{\phi}] \leftarrow L$ \tcp*[l]{choose the right most subgoal}
    $\Pi_1 \seqA \msgB{op}{\overline{y}}{\phi} \seqA \Pi_2 \leftarrow \Pi$  \tcp*[l]{divide the symbolic plan according to the select subgoal}
    $E \leftarrow \emptyset$\;
    \ForEach{$(\Gamma', \Pi') \in \S{BackwardSynthesize}(\Delta, L, \Gamma, \Pi_1, \msgB{op}{\overline{y}}{\phi}, \Pi_2)$}{
        $E \leftarrow E \cup \S{Loop}(\Gamma, \Pi)$\;
    }
    \Return{$E$}\;
}
    \caption{Synthesis}
    \label{algo:top-syn}
  \end{algorithm}

\begin{algorithm}[t!]
\Params{Effect operator type context $\Delta$, record subgoals $L$, type context $\Gamma$, history automata $\Pi_1$, current goal $\msgB{op_2}{\overline{y}}{\phi}$, future automata $\Pi_2$}
\Output{A set of pair $(\Gamma, \Pi)$}
% $S \leftarrow \emptyset$ \tcp*[l]{a set record all subgoals that are solved in this backward synthesis procedure}
% $G \leftarrow \S{ForwardSynthesis}(\Delta, L, \Gamma, \Pi_1, \msgB{op_2}{\overline{x}}{\phi} \seqA \Pi_2)$ \tcp*[l]{do forward synthesis}
\ForEach{$(\Gamma, \Pi_1, \eff{op_2}
(\overline{x'}) \seqA \Pi_2) \in \S{ForwardSynthesis}(\Delta, L, \Gamma, \Pi_1, \msg{op_2}{\phi} \seqA \Pi_2)$}{
    \tcc{forward synthesis guarantees that $\eff{op_2}$ is solved. }
    \tcc{If $\eff{op_2}$ can be directly generated by environment, backward synthesis ends}
    \lIf{$\S{Gen}(\eff{op_2})$}{
        \Return{$\{(\Gamma, \Pi_1 \seqA \eff{op_2}
(\overline{x'}) \seqA \Pi_2)\}$}
    }
    $S \leftarrow \emptyset$ \tcp*[l]{A set of all synthesized results}
    \tcc{The handling procedure of message $\eff{op_1}$ will send new message $\eff{op_2}$}
    \ForEach{$(\eff{op_1}, \overline{z{:}b}\garr\overline{y{:}t}\sarr\rg{\Pi_h}{\msg{op_1}{\phi_1}}{\Pi_f}) \in \Delta$ where $\eff{op_2} \in \S{SendEvent}(\Pi_f)$}{
        $\Pi_{f1} \seqA \msgB{op_2}{\overline{x}}{\phi_2} \seqA \Pi_{f2} \leftarrow \Pi_f$\;
        $\Gamma \leftarrow \Gamma, \overline{z{:}\rawnuot{b}{\top}}, \overline{y{:}t}$ \tcp*[l]{add parameters of $\eff{op_1}$ into type context}
        \tcc{align automata before $\eff{op_2}$}
        \ForEach{$\Pi_h' \seqA \msg{op_1}{\phi_1'} \seqA \Pi_{f1}' \in \S{merge}(\Pi_1,\; \Pi_h \seqA \msg{op_1}{\phi_1} \seqA \Pi_{f1})$}{
            \tcc{align automata after $\eff{op_2}$}
            \ForEach{$\Pi_{f2}' \in \S{merge}(\Pi_2,\; \Pi_{f2})$}{
                \tcc{Refine $\Gamma$ (i.e., $\overline{z}$ and $\overline{y}$) to make the total automata is nonempty.}
                \Match{$\S{Abduction}(\Gamma, \Pi_h' \seqA \msg{op_1}{\overline{x = y} \land \phi_1'} \seqA \Pi_{f1}' \seqA \msgB{op_2}{x}{\overline{x = x'} \land \phi_2} \seqA \Pi_{f2}')$}{
                    \lCase{$\Code{None}$}{ \Continue }
                    \Case{$\Code{Some}\;\Gamma'$}{
                        \tcc{now $\eff{op_1}$ is solved.}
                        $S \leftarrow S \cup \S{BackwardSynthesis}(\Delta, L, \Gamma', \Pi_h', \eff{op_1}(\overline{y}), \Pi_{f1}' \seqA \eff{op_2}(\overline{x'}) \seqA \Pi_{f2}')$\;
                    }
                }
            }
        }
    }
    \Return{$S$}\;
}
    \caption{Backward Synthesis}
    \label{algo:top-syn}
\end{algorithm}

\begin{algorithm}[t!]
\Params{Effect operator type context $\Delta$, record subgoals $L$, type context $\Gamma$, history automata $\Pi_1$, future automata $\Pi_2$}
\Output{A set of pair $(\Gamma, \Pi_1, \Pi_2)$ guarantees that no $\msgB{op}{\overline{y}}{\phi}$ in $\Pi_2$ that is not record by $L$.}
% $S \leftarrow \emptyset$ \tcp*[l]{a set record all subgoals that are solved in this backward synthesis procedure}
\tcc{select a $\msgB{op}{\overline{y}}{\phi}$ in $\Pi_2$.}
\uIf{$\exists\; \Pi_{21}\; \msg{op}{\phi}\; \Pi_{22}.\; \Pi_1 = \Pi_{21} \seqA \msg{op}{\phi} \seqA \Pi_{22} \land \msg{op}{\phi} \not\in Q$}{
    $S \leftarrow \emptyset$\;
    \ForEach{$(\eff{op}, \overline{z{:}b}\garr\overline{y{:}t}\sarr\rg{\Pi_h}{\msg{op}{\phi'}}{\Pi_f}) \in \Delta$}{
        $\Gamma \leftarrow \Gamma, \overline{z{:}\rawnuot{b}{\top}}, \overline{y{:}t}$\;
        \tcc{align automata before $\eff{op}$}
        \ForEach{$\Pi_1' \seqA \Pi_{21}' \in \S{merge}(\Pi_1 \seqA \Pi_{21},\; \Pi_h)$}{
            \tcc{align automata before $\eff{op}$}
            \ForEach{$\Pi_{22}' \in \S{merge}(\Pi_{22},\; \Pi_{f})$}{
                \tcc{Refine $\Gamma$ (i.e., $\overline{z}$ and $\overline{y}$) to make the total automata is nonempty.}
                \Match{$\S{Abduction}(\Gamma, \Pi_1' \seqA \Pi_{21}' \seqA \msgB{op}{\overline{x}}{\overline{x = y} \land \phi \land \phi'} \seqA \Pi_{22}')$}{
                    \lCase{$\Code{None}$}{ \Continue }
                    \Case{$\Code{Some}\;\Gamma'$}{
                        \tcc{now $\eff{op}$ is solved.}
                        $S \leftarrow S \cup \S{ForwardSynthesis}(\Delta, L, \Gamma', \Pi_1', \Pi_{21}' \seqA \eff{op}(\overline{y}) \seqA \Pi_{22}')$\;
                    }
                }
            }
        }
    }
    \Return{$S$}\;
}
\Else{\Return{$\{(\Gamma, \Pi_1, \Pi_2)\}$}}
    \caption{Forward Synthesis}
    \label{algo:top-syn}
\end{algorithm}



\paragraph{Symbolic Plan} As illustrated in \autoref{fig:plan-syntax}, the symbolic plan is similar to symbolic automata but presented in a sequential format. In addition to the symbolic event $\msgB{op}{\overline{x}}{\phi}$, we introduce \emph{symbolic applications} $\acSe{op}{\overline{x}}$, which represent the application of effect operators. Unlike concrete events $\eff{op}(\overline{c})$, symbolic applications allow variables to be passed as arguments. A symbolic application can be viewed as a special type of symbolic event $\msgB{op}{\overline{y}}{\bigwedge \overline{y = x}}$, where all arguments are properly instantiated. We also introduce a special Kleene star operator $A^\acStar$, which allows us to encapsulate arbitrarily complex symbolic finite automata (SFA) $A$ without unrolling them into a sequential format. This enables lazy unrolling of the Kleene star as needed, improving efficiency.  Finally, unlike SFAs, symbolic plans are free of union operators. As a result, an SFA can be normalized (via the auxiliary function $\normPlan$) into multiple symbolic plans. Importantly, since we use the lazy star $A^\acStar$, the normalization function $\normPlan$ always produces a finite number of plans.

\paragraph{Synthesis Judgements} The planing is just try to convert all $\msgB{op}{\overline{x}}{\phi}$ into $\acSe{op}{\overline{v}}$. There are three classes of synthesis judgements:
\begin{enumerate}
    \item Global Synthesis $\Gamma \vdash A \terminfer e$ and $\Gamma \vdash \Pi \terminfer e$.
    \item Forward Synthesis $\Gamma \vdash \planPair{\Pi_1}{\Pi_2} \inferFW \planPair{\Pi_1'}{\Pi_2'}$, needs to guarantee that there is not $\msgB{op}{\overline{x}}{\phi}$ in $\Pi_2'$.
    \item Backward Synthesis $\Gamma \vdash \planTri{\Pi_1}{\acSe{op}{\overline{x}}}{\Pi_2} \inferBW \Pi'$, needs to guarantee that all $\acSe{op}{\overline{x}}$ in $\Pi'$ are consistent with $\Delta$.
\end{enumerate}

\newcommand{\aaStar}{\anyA^\acStar}
Here is a simple intuition how forward/backward works:
\begin{align*}
    \effRed{WriteReq^5} \to \effRed{PutReq^4} \to \effRed{PutRsp^2} &\to \effRed{WriteRsp^1}
    \\&\searrow \effRed{Commit^3}
\end{align*}\noindent
Assume we start from $\vdash \aaStar \seqA \eff{WriteReq} \seqA \aaStar$. 
\begin{enumerate}
    \item The rule \textsc{SynFocus} makes us focus on the symbolic event $\eff{WriteReq}$. We first do forward synthesis $\vdash \planPair{\aaStar}{\eff{WriteReq} \seqA \aaStar}$ to make sure all messages trigger by $\eff{WriteReq}$ are refined; then we do backward synthesis from $\eff{WriteReq}$.
    \item $\eff{WriteReq}$ doesn't trigger new events. The forward synthesis result is $\vdash \planPair{\aaStar}{\eff{WriteReq} \seqA \aaStar} \inferFW \vdash \planPair{\aaStar}{\effRed{WriteReq} \seqA \aaStar}$.
    \item We do backward synthesis $\vdash \planTri{\aaStar}{\effRed{WriteReq}}{\aaStar}$. As we know $\eff{PutRsp}$ trigger it. We have $\vdash \planTri{\aaStar \seqA \effRed{PutRsp} \seqA \aaStar}{\effRed{WriteReq}}{\aaStar \seqA \eff{Commit} \seqA \aaStar}$. The invariant of ``no symbolic event after $\planBar$'' is broken, thus we apply forward synthesis $\vdash \planPair{\aaStar \seqA \effRed{PutRsp} \seqA \aaStar}{\effRed{WriteReq} \seqA \aaStar \seqA \eff{Commit} \seqA \aaStar}$ again.
    \item $\eff{Commit}$ doesn't trigger new events. The forward synthesis result is 
    \begin{align*}
        \vdash \planPair{\aaStar \seqA \effRed{PutRsp} \seqA \aaStar}{\effRed{WriteReq} \seqA \aaStar \seqA \eff{Commit} \seqA \aaStar} \inferFW \planPair{\aaStar \seqA \effRed{PutRsp} \seqA \aaStar}{\effRed{WriteReq} \seqA \aaStar \seqA \effRed{Commit} \seqA \aaStar}
    \end{align*}
    \item We continue backward synthesis, similarly, we have
    \begin{align*}
        &\vdash \planTri{\aaStar}{\effRed{PutRsp}}{\aaStar \seqA \effRed{WriteReq} \seqA \aaStar \seqA \effRed{Commit} \seqA \aaStar} 
        \\&\inferBW \planTri{\aaStar}{\effRed{PutReq}}{\aaStar \seqA \effRed{PutRsp} \seqA \aaStar \seqA \effRed{WriteReq} \seqA \aaStar \seqA \effRed{Commit} \seqA \aaStar}
        \\&\inferBW \planTri{\aaStar}{\effRed{WriteReq}}{\aaStar \seqA \effRed{PutReq} \seqA \aaStar \seqA \effRed{PutRsp} \seqA \aaStar \seqA \effRed{WriteReq} \seqA \aaStar \seqA \effRed{Commit} \seqA \aaStar}
    \end{align*}
    \item Finally, we $\effRed{WriteReq}$ can be generated by environment directly, then backward synthesis ends.
\end{enumerate}

{\footnotesize\mprooftr{46mm}{
$\text{no $\msgB{op}{\overline{x}}{\phi}$ in } \Pi$ \quad
$\Pi' = \Pi[A^\acStar \mapsto \epsilon]$
}
{SynStar}
{$\Gamma \vdash \Pi \terminfer \termDerive(\Gamma, \Pi')$}}
\quad 
{\footnotesize\mprooftr{46mm}{
$\Pi_i \in \normPlan(A)$ \quad
$\Gamma \vdash \Pi_i \terminfer e_i$
}
{SynUnion}
{$\Gamma \vdash A \terminfer \bigoplus e_i$}}
\\ \ \\ \ 

{\footnotesize\mprooftr{65mm}{
$\text{no $\msgB{op}{\overline{x}}{\phi}$ in } \Pi_2$  \quad
% $\Gamma \vdash \planPair{ \Pi_1 }{ \msgB{op}{\overline{x}}{\phi} \seqA \Pi_2 } \inferFW \Gamma' \vdash \planPair{ \Pi_1' }{ \acSe{op}{\overline{y}} \seqA \Pi_2' }$ \\
$\Gamma \vdash \planTri{ \Pi_1 }{ \acSe{op}{\overline{y}} }{ \Pi_2 } \inferBW \Gamma' \vdash \Pi'$ \\
$\Gamma' \vdash \Pi' \terminfer e$
}
{SynFocus}
{$\Gamma \vdash \Pi_1 \seqA \msgB{op}{x}{\phi} \seqA \Pi_2 \terminfer e$}}
\\ \ \\ \ 

{\footnotesize\mprooftr{80mm}{
% $\Pi_{21} \seqA \msgB{op}{\overline{x}}{\phi} \seqA \Pi_{22} \in \normPlan(\Pi_2)$\quad
$\overline{y{:}b} \sarr [H]\msgB{op}{x}{\phi'}[F] \in \Delta(\eff{op})$ \\
$\Pi_1' \seqA \Pi_{21}' \in \normPlan(\Pi_1 \seqA \Pi_{21} \;\land\; H)$ \quad
$\Pi_{22}' \in \normPlan(\Pi_{22} \;\land\; F)$ \\
$\Gamma \infervdash{\Pi_1' \seqA \Pi_{21}' \seqA \msgB{op}{x}{\phi \land \phi'} \seqA \Pi_{22}'} \overline{y{:}b} \effinfer \overline{y{:}\rawnuot{b}{\phi_y}}$  \\
$\Gamma, \overline{y{:}\rawnuot{b}{\phi_y}} \vdash \planPair{ \Pi_1' }{ \Pi_{21}' \seqA \acSe{op}{\overline{y}} \seqA \Pi_{22}' } \inferFW \Gamma' \vdash \planPair{ \Pi_3 }{ \Pi_4 }$ 
}
{FWAct}
{$\Gamma \vdash \planPair{ \Pi_1 }{ \Pi_{21} \seqA \msgB{op}{\overline{x}}{\phi} \seqA \Pi_{22} } \inferFW \Gamma' \vdash \planPair{\Pi_3}{\Pi_4}$}}
\quad
{\footnotesize\mprooftr{40mm}{
$\text{no $\msgB{op}{\overline{x}}{\phi}$ in } \Pi_2$
}
{FWEnd}
{$\Gamma \vdash \planPair{\Pi_1}{\Pi_2} \inferFW \Gamma \vdash \planPair{\Pi_1}{\Pi_2}$}}
\\ \ \\ \ 

{\footnotesize\mprooftr{110mm}{
$\text{no $\msgB{op}{\overline{x}}{\phi}$ in } \Pi_2$ \quad
$\overline{y{:}b} \sarr [H]\msgB{op_1}{\overline{x}}{\phi_1}[F_1 \seqA \msgB{op_2}{\overline{z}}{\phi_2} \seqA F_2 ] \in \Delta(\eff{op_1})$ \quad 
$\Pi_{11}'\seqA\msgB{op_1}{\overline{x}}{\phi_1'}\seqA\Pi_{12}' \in \normPlan(\Pi_1 \land H\seqA\msgB{op_1}{\overline{x}}{\phi_1}\seqA F_1)$ \quad
$\Pi_{2}' \in \normPlan(\Pi_2 \land F_2)$ \\
$\Gamma, \infervdash{ \Pi_{11}'\seqA\msgB{op_1}{\overline{x}}{\phi_1'}\seqA\Pi_{12}' \seqA \msgB{op_2}{\overline{z}}{\overline{z = z'} \land \phi_2} \seqA \Pi_{2}'} \overline{y{:}b} \effinfer \overline{y{:}\rawnuot{b}{\phi_y}}$ \\
$\Gamma, \overline{y{:}\rawnuot{b}{\phi_y}} \vdash \planTri{ \Pi_{11}'}{\acSe{op_1}{\overline{y}} }{ \Pi_{12}' \seqA \acSe{op_2}{\overline{z'}} \seqA \Pi_{2}' } \inferBW \Gamma' \vdash \Gamma_3$ 
% \\
% $\Gamma'' \vdash \planTri{ \Pi_3}{\acSe{op_1}{\overline{y}} }{ \Pi_4 } \inferBW \Gamma''' \vdash \Pi_5$
}
{BWAct}
{$\Gamma \vdash \planTri{ \Pi_1 }{ \acSe{op_2}{\overline{z'}} }{ \Pi_2 } \inferBW \Gamma' \vdash \Gamma_3$}}
\\ \ \\ \ 

{\footnotesize\mprooftr{55mm}{
$\text{no $\msgB{op}{\overline{x}}{\phi}$ in } \Pi_2$ \quad
$\S{Gen}(\eff{op})$}
{BWEnd}
{$\Gamma \vdash \planTri{\Pi_1}{ \acSe{op}{\overline{x}} }{\Pi_2} \inferBW \Gamma \vdash \Pi_1 \seqA \acSe{op}{\overline{x}} \seqA \Pi_2$}}
\quad
{\footnotesize\mprooftr{55mm}{
$\text{has $\msgB{op}{\overline{x}}{\phi}$ in } \Pi_2$ \quad
$\Gamma \vdash \planPair{\Pi_1}{ \acSe{op}{\overline{x}} \seqA \Pi_2} \inferFW \Gamma' \vdash \planPair{\Pi_1'}{ \acSe{op}{\overline{x}} \seqA \Pi_2'}$\quad
$\Gamma' \vdash \planTri{\Pi_1'}{ \acSe{op}{\overline{x}} }{\Pi_2'} \inferBW \Gamma'' \vdash \Pi'$
}
{BWFW}
{$\Gamma \vdash \planTri{\Pi_1}{ \acSe{op}{\overline{x}} }{\Pi_2} \inferBW \Gamma'' \vdash \Pi'$}}

\paragraph{Term Derivation}

\begin{algorithm}[t!]
    \lProcedure{$\termDerive(\Gamma, \Pi) := $}{
        \Return{$\S{DerivePlan}(\Pi, \Gamma, \emptyset)$}
    }
    \Procedure{$\S{DerivePlan}(\Pi, \Gamma_{\Pi}, \Gamma_{e}) := $}{
    \Match{$\Pi$}{
      \lCase{$\epsilon$}{ \Return{$()$} }
      \Case{$\acSe{op}{\overline{x}} \seqA \Pi'$}{
        $\overline{x'}, \overline{x''} \leftarrow \overline{x'}$ where $ \overline{x''} \subseteq \S{Dom}(\Gamma_{e})$\;
        $\Gamma_{\Pi}' \leftarrow \Gamma_{\Pi} \setminus \overline{x'}$\;
        $\Gamma_{e} \infervdash{\exists \Gamma_\Pi'. \Pi} \overline{x'} \effinfer \overline{x'{:}\rawnuot{b}{\phi}}$\;
        $\Gamma_{e}' \leftarrow \Gamma_{e}, \overline{x'{:}\rawnuot{b}{\phi}}$\;
        $e \leftarrow \S{DerivePlan}( \Pi', \Gamma_{\Pi}', \Gamma_{e}')$\;
        \uIf{$\S{Gen}(\eff{op})$}{
            \lForEach{$x_i{:}\rawnuot{b}{\phi_i}$}{
                $e_i \leftarrow \S{DeriveArg}(\Gamma_e, \rawnuot{b}{\phi_i})$
            }
            \Return{$\zlet{\overline{x_i}}{\overline{e_i}}{\zgen{op}{\overline{e_i}  \; \overline{x''}}{e}}$}\;
        }
        \Else{
            $\phi \leftarrow \top; \quad \overline{y} \leftarrow \emptyset$\;
            \lForEach{$x_i'{:}\rawnuot{b'}{\phi_i'}$}{
                $\phi \leftarrow \phi \land \phi_i'[\vnu \mapsto x_i']$
            }
            \ForEach{$x_i'' \in \overline{x''}$}{
                $y_i \leftarrow  \S{freshVar}()$\;
                $\overline{y} \leftarrow  \overline{y}, y_i; \quad \phi \leftarrow \phi \land y_i = x_i''$\;
            }
            \Return{$\zobs{\overline{x'}\; \overline{y}}{op}{\zassert{\phi}{e}}$}\;
        }
      }
    }
    }
    \caption{Term Derivation}
    \label{algo:term-derive}
  \end{algorithm}

\paragraph{Plan Normalization}

\begin{algorithm}[t!]
    \Procedure{$\normPlan(A) := $}{
    \Match{$A$}{
      \lCase{$\emptyset$}{ \Return{$\emptyset$} }
      \lCase{$\epsilon$}{ \Return{$\{\epsilon\}$} }
      \lCase{$\msgB{op}{\overline{x}}{\phi}$}{ \Return{$\{\msgB{op}{\overline{x}}{\phi}\}$} }
      \lCase{$A_1 \lorA A_2$}{ 
      \Return{$\{ \Pi_1 \seqA \Pi_2 ~|~ \Pi_1 \in \normPlan(A_1) \land \Pi_2 \in \normPlan(A_2) \}$}}
      \lCase{$A^*$}{ \Return{$\{ A^\acStar \}$}}
      \lCase{$\anyA$}{ \Return{$\{ \msgB{op}{\overline{x}}{\top} ~|~ \text{for all } \eff{op} \}$}}
    }
    }
    \caption{Plan Normalization}
    \label{algo:norm}
  \end{algorithm}